\documentclass[11pt]{article}
\usepackage[T1]{fontenc}
\usepackage{lmodern}
\usepackage[utf8]{inputenc}
\usepackage[margin=1in]{geometry}
\usepackage{amsmath,amssymb,amsthm,mathtools}
\usepackage[hidelinks]{hyperref}

\newtheorem{definition}{Definition}
\newtheorem{theorem}{Theorem}
\newtheorem{proposition}{Proposition}
\newtheorem{remark}{Remark}
\newcommand{\Sub}{\mathcal S}
\newcommand{\Lie}{\mathcal L}
\newcommand{\Char}{\operatorname{Char}}
\newcommand{\Ann}{\operatorname{Ann}}
\newcommand{\supp}{\operatorname{supp}}

\title{A Proposal-Only Source-Local Transport Characteristic\\
\large P6-T02 registered draft/control construction}
\author{The AEther-Flow Research Project}
\date{2026-07-26}

\begin{document}
\maketitle

\begin{abstract}
This registered draft/control manuscript executes v21 work item P6-T02.  It
adds one materially new but proposal-only source extension to the canonical
smooth four-dimensional source arena: a scalar amplitude field, a nowhere-zero
source vector field, and the local transport--reaction equation
\(\Lie_V a+\gamma a^3=0\).  The equation exactly reproduces the P5 cubic flow
along integral curves.  Its linearized principal polynomial
\(P_x(\xi)=\xi(V_x)\), characteristic covectors, hyperbolicity half-space,
dual propagation ray, and forward reachability preorder are calculated from
the equation without a target metric or GR null cone.  An exact response
formula, a nontrivial flow-box example, an \(O(\varepsilon)\) upwind
consistency bound, and zero-field, closed-orbit, competing-generator, and
multi-sector controls are supplied.  The result is a constructed formal
source-extension candidate, not a physical causal structure: the current
ontology does not derive or select \(V\), its orientation, normalization,
global chronology, operational signal meaning, or universality.  No
effective conformal class, metric scale, matter coupling, Einstein equation,
ontology adoption, or benchmark promotion follows.
\end{abstract}

\section{Authority and status boundary}

The canonical source architecture carries a smooth four-dimensional manifold
\(\Sub\) as explicit primitive debt and leaves \(\Phi_{\mathrm{src}}\) as an
unresolved source-order or evolution slot.  It does not identify \(\Sub\)
with physical spacetime and does not adopt a generator, source field,
localization law, causal order, metric, or source-to-target bridge.

P5 supplies only the proposal-only scalar flow
\[
 \phi_s(z)=\frac{z}{\sqrt{1+2\gamma s z^2}},
 \qquad
 D\phi_s(z)=(1+2\gamma s z^2)^{-3/2},
 \qquad \gamma>0,\ s\geq0.
\]
P6-T01 defines a metric-free characteristic target and proves that these
unchanged diagonal data do not select a unique region-resolved influence
relation.  The present packet does not reopen that inference.  It introduces
new source-local law data and derives their consequences conditionally.

Every new object below has status
\texttt{draft/control}, \texttt{proposal-only}, and
\texttt{source-extension data}.  In particular, the vector field \(V\) is
not a target-side congruence, preferred physical frame, ordinary wind,
physical clock, or canonical-ontology candidate.

\section{Candidate and characteristic calculation}

\begin{definition}[Source-local transport candidate]
\label{def:candidate}
Let \(\Sub\) be a smooth four-dimensional source manifold.  On an open set
\(U\subseteq\Sub\), propose:
\begin{enumerate}
  \item a scalar source field \(a\in C^\infty(U,\mathbb R)\);
  \item a smooth nowhere-zero source vector field \(V\in\Gamma(TU)\);
  \item a source-side orientation selecting the positive ray of \(V\); and
  \item a constant \(\gamma>0\).
\end{enumerate}
Define the local equation
\[
 \mathcal E_V[a]
 :=\Lie_Va+\gamma a^3
 =V^A\partial_Aa+\gamma a^3=0.
\tag{1}\label{eq:candidate}
\]
These are proposed source data.  No target coordinate, metric, connection,
volume form, null cone, matter equation, or benchmark solution enters
\eqref{eq:candidate}.
\end{definition}

\begin{theorem}[P6T02-THM-FLOW-BOX-RECOVERY]
\label{thm:flowbox}
Let \(\chi_s\) be a positive integral curve of \(V\), with
\(\chi_0=x_0\), on an interval containing \(s\geq0\).  A solution of
\eqref{eq:candidate} obeys
\[
 a(\chi_s)
 =\frac{a(x_0)}{\sqrt{1+2\gamma s\,a(x_0)^2}}.
\tag{2}\label{eq:solution}
\]
Thus the restriction of the proposed local equation to every integral curve
is exactly the P5 cubic forward flow.  In a flow box
\((s,y^1,y^2,y^3)\) with \(V=\partial_s\), initial data \(a_0(y)\) on
\(s=0\) give
\[
 a(s,y)=\frac{a_0(y)}
 {\sqrt{1+2\gamma s\,a_0(y)^2}}.
\tag{3}\label{eq:flowboxsolution}
\]
\end{theorem}

\begin{proof}
Along an integral curve,
\(\frac{d}{ds}a(\chi_s)=\Lie_Va=-\gamma a^3\).  Separation and the
initial condition give \eqref{eq:solution}.  The flow-box expression follows
pointwise in the transverse labels \(y\).
\end{proof}

\begin{theorem}[P6T02-THM-PRINCIPAL-RAY]
\label{thm:principal}
Linearization about any smooth background \(\bar a\) satisfying
\eqref{eq:candidate} gives
\[
 \Lie_Vu+3\gamma\bar a^2u=0.
\tag{4}\label{eq:linearized}
\]
Its scalar principal symbol and principal polynomial are
\[
 \sigma_x(\xi)u=\xi(V_x)u,
 \qquad
 P_x(\xi)=\xi(V_x).
\tag{5}\label{eq:principal}
\]
Consequently,
\[
 \Char_x^*=\{\xi\in T_x^*\Sub\setminus\{0\}:\xi(V_x)=0\}
 =\Ann(V_x)\setminus\{0\}.
\tag{6}\label{eq:char}
\]
For any covector \(h\) with \(h(V_x)>0\), the degree-one polynomial
\(P_x\) is hyperbolic with respect to \(h\), its selected hyperbolicity
 component is
\[
 \Gamma_x=\{\xi:\xi(V_x)>0\},
\tag{7}\label{eq:halfspace}
\]
and the closed dual cone under the natural tangent--cotangent pairing is the
single ray
\[
 \Gamma_x^\circ=\{rV_x:r\geq0\}.
\tag{8}\label{eq:dualray}
\]
The base projection of the principal Hamiltonian direction is \(V\).
Therefore the equation supplies one formal oriented source propagation ray,
not a Lorentzian cone or metric.
\end{theorem}

\begin{proof}
The cubic term is zeroth order, so only \(\Lie_Vu\) contributes to the
principal symbol.  For fixed \(\xi\),
\[
 \lambda\longmapsto P_x(\xi+\lambda h)
 =\xi(V_x)+\lambda h(V_x)
\]
has one real root because \(h(V_x)\neq0\).  The selected component is the
positive half-space \eqref{eq:halfspace}.  If a tangent vector pairs
nonnegatively with every covector in that half-space, separation of a vector
from the line spanned by \(V_x\) shows that it must be a nonnegative multiple
of \(V_x\), proving \eqref{eq:dualray}.  Finally,
\(\partial P/\partial\xi=V\), which gives the base propagation direction.
\end{proof}

\begin{proposition}[P6T02-PROP-RESPONSE-SUPPORT]
\label{prop:response}
In a flow box, the solution of the linearized equation with initial
perturbation \(u_0(y)\) is
\[
 u(s,y)=u_0(y)
 \bigl(1+2\gamma s\,a_0(y)^2\bigr)^{-3/2}.
\tag{9}\label{eq:response}
\]
Hence perturbation support is transported only along the positive \(V\)-flow:
\[
 \supp u(s,\cdot)
 =\chi_s\bigl(\supp u_0\bigr)
\]
whenever the flow map is defined and the nonzero multiplier is interpreted
on corresponding flow lines.  The relation
\[
 x\preceq_V y
 \quad\Longleftrightarrow\quad
 y=\chi_s(x)\ \text{for some }s\geq0
\tag{10}\label{eq:preorder}
\]
is therefore the candidate's formal forward influence preorder.
\end{proposition}

\begin{proof}
Differentiate \eqref{eq:solution} with respect to its initial value.  The
multiplier is exactly the P5 tangent response and is strictly positive for
finite \(s\geq0\), so it neither creates nor cancels support between distinct
flow lines.  Flow composition makes \(\preceq_V\) reflexive and transitive.
\end{proof}

\section{Covariance, normalization, and sector count}

The construction is source-diffeomorphism covariant when \(a\) is transported
as a scalar and \(V\) as a vector field.  The characteristic set and ray are
also invariant under positive rescaling \(V\mapsto fV\), \(f>0\), because
\(P\mapsto fP\).  The full equation is multiplied by \(f\) only if the
reaction coefficient is simultaneously treated as the source function
\(\gamma\mapsto f\gamma\).  Thus the unparameterized oriented line field
determines the formal characteristic ray, but the exact P5 parameterization
still requires an independent normalization or scale law.  No such law is
derived here.

The scalar candidate has one propagation ray at each regular point.  A
proposal-only two-sector direct sum with vector fields \(V\) and \(W\) has
principal determinant
\[
 P^{(2)}_x(\xi)=\xi(V_x)\xi(W_x),
\tag{11}\label{eq:twosector}
\]
and two characteristic hyperplanes or propagation rays when \(V_x\) and
\(W_x\) are noncollinear.  The present sources contain no universality rule
that selects a common ray across such sectors.

\section{Nontrivial and pathological controls}

\paragraph{Nontrivial flow-box witness.}
On \(\Sub=\mathbb R^4\) with coordinates \((s,y^1,y^2,y^3)\), choose
\(V=\partial_s\) and
\(a_0(y)=A\exp[-((y^1)^2+(y^2)^2+(y^3)^2)]\), \(A\neq0\).
Equation~\eqref{eq:flowboxsolution} is a smooth, nonconstant exact solution.
Its characteristic covectors satisfy \(\xi_s=0\), and its formal forward
propagation ray is the positive \(s\)-direction.

\paragraph{Zero-field degeneracy.}
If \(V_x=0\), then \(P_x\equiv0\) on \(T_x^*\Sub\).  There is no
hyperbolicity direction and no nondegenerate propagation ray at \(x\).
Nowhere-vanishing is therefore a genuine candidate hypothesis, not a
conclusion.  Some global four-manifolds do not admit such a vector field
without additional topology assumptions.

\paragraph{Closed-orbit chronology failure.}
On \(S^1\times\mathbb R^3\), take \(V=\partial_\theta\).  Distinct points on
the same periodic orbit satisfy both \(x\preceq_Vy\) and
\(y\preceq_Vx\).  The relation is a preorder but not an antisymmetric causal
order.  A no-closed-orbit or source-time function condition is an additional
global hypothesis.

\paragraph{Competing-generator nonselection.}
On \(\mathbb R^4\), both
\(V=\partial_s\) and \(W=\partial_s+\kappa\partial_{y^1}\),
\(\kappa\neq0\), recover the same P5 equation on fields independent of the
transverse coordinates, but their characteristic hyperplanes and propagation
rays differ.  The unchanged P5 data and current ontology do not select
between them.  This is not a global impossibility theorem; it precisely
locates the missing selector.

\section{Finite approximation and error}

In the flow-box witness, approximate \(\partial_sa\) by the backward
source-oriented difference
\[
 D^-_\varepsilon a(s,y)
 =\frac{a(s,y)-a(s-\varepsilon,y)}{\varepsilon}.
\tag{12}\label{eq:upwind}
\]
For \(a\in C^2\) on the sampled segment, Taylor's theorem gives
\[
 \left|D^-_\varepsilon a-\partial_sa\right|
 \leq
 \frac{\varepsilon}{2}
 \sup_{\rho\in[s-\varepsilon,s]}
 |\partial_s^2a(\rho,y)|.
\tag{13}\label{eq:error}
\]
Thus the local source-operator reconstruction is first-order accurate,
\(O(\varepsilon)\).  This is not an effective-metric reconstruction error.
No source-to-target reconstruction map is available, so a target conformal
class or metric error is undefined rather than silently set to zero.

\section{Disposition}

The decisive result is \texttt{constructed\_candidate}.  The proposed local
equation yields, by calculation, one principal polynomial, characteristic
hyperplane, hyperbolicity half-space, dual ray, exact response-support law,
and formal forward preorder.  This is materially new source-local law content
and does not repeat the P6-T01 two-site inference.

Physical Gate B remains not ready.  The current ontology does not derive or
select the scalar field, \(V\), its positive orientation, normalization,
global nonvanishing, absence of closed orbits, signal-sector universality, or
operational meaning.  The one-ray cone is also too degenerate by itself to be
called a four-dimensional Lorentzian conformal structure.  P6-T03 may test
that conformal insufficiency after the governed checkpoint.

\section{Explicit nonclaims}

This packet does not edit canonical ontology; adopt or reject a source law;
identify \(\Sub\) with physical spacetime; identify the integral-curve
parameter with physical time; establish physical locality, causality,
chronology, a clock, or a signal; select \(V\) uniquely; import a target
atlas, GR null cone, proper time, detector, metric, or benchmark solution;
construct an unscoped \(g_{\mathrm{eff}}\); alter the historical scoped
\(g_{\mathrm{eff}}\) object; couple matter; derive Einstein equations;
promote the exact-GR benchmark; issue a Gate Chair verdict; prove a global
no-go theorem; declare future source extensions impossible; publish; push; or
claim a completed derivation.

\end{document}
